\documentclass[12pt]{article}
\usepackage[T1]{fontenc}
\usepackage[utf8]{inputenc}
\usepackage[italian]{babel}
\usepackage{amsfonts, amssymb, amsmath}

\title{Riassunto Didattico: Logica di Coalizione e Quantificazione}
\author{Studente: [Il Tuo Nome]}
\date{\today}

\begin{document}

\maketitle

% PIANIFICAZIONE DELLA PRESENTAZIONE (20 MINUTI TOTALI)
% 00:00 - 03:00 -> Introduzione e Motivazione (Sezione 1)
% 03:00 - 08:00 -> Fondamenta: Modelli di Gioco ed Effettività (Sezione 2 e 3)
% 08:00 - 10:00 -> La Coalition Logic (CL) e i suoi limiti (Sezione 4)
% 10:00 - 15:00 -> La Svolta: Quantified Coalition Logic (QCL) (Sezione 5 e 6)
% 15:00 - 18:00 -> Risultati Teorici: Complessità e Completezza (Sezione 7)
% 18:00 - 20:00 -> Esempi Pratici e Conclusioni (Sezione 8)

\section{Introduzione: Dai Modelli Kripke ai Frame di Gioco}
% SUGGERIMENTO DIDATTICO: Qui devi spiegare perché la logica modale standard non basta.
% Passiamo da un solo agente a un'interazione simultanea tra molti [6, 7].

\section{Il Concetto di Potere Coalizionale}
% SUGGERIMENTO DIDATTICO: Spiega la funzione di effettività.
% Cosa significa "forzare" un esito (E-effectivity) [8, 9]?

\section{Coalition Logic (CL): Sintassi e Semantica}
% SUGGERIMENTO DIDATTICO: Introduciamo l'operatore [[C]]phi [2, 10].
% È qui che parleremo del Teorema 3.2 (Caratterizzazione) che abbiamo visto prima [11].

\section{Assiomatizzazione e Limiti della CL}
% SUGGERIMENTO DIDATTICO: Elenca gli assiomi base (superadditività, ecc.) [12].
% Introduci il problema della "succinctness": perché enumerare tutte le coalizioni è inefficiente [13, 14]?

\section{Quantified Coalition Logic (QCL): Predicati di Coalizione}
% SUGGERIMENTO DIDATTICO: Questa è la parte nuova della seconda fonte.
% Spiega come usiamo sottoinsiemi, sovrainsiemi e cardinalità per descrivere gruppi (subseteq, geq) [15, 16].

\section{Operatori Quantificati: Esistenziale e Universale}
% SUGGERIMENTO DIDATTICO: Introduciamo <P>phi e [P]phi [17, 18].
% Mostra come questi operatori rendono la logica esponenzialmente più concisa [5].

\section{Metateoria: Complessità Computazionale}
% SUGGERIMENTO DIDATTICO: Un punto fondamentale per il tuo pubblico sarà che,
% nonostante la QCL sia più potente espressivamente, rimane PSPACE-completa come la CL [19, 20].

\section{Applicazioni: Meccanismi di Voto e Scelte Sociali}
% SUGGERIMENTO DIDATTICO: Usa l'esempio del voto a maggioranza per mostrare
% la bellezza della QCL rispetto alla CL [21, 22].

\section{Conclusioni e Riflessioni Personali}
% SUGGERIMENTO DIDATTICO: Riassumi il legame tra teoria dei giochi e logica formale.

\end{document}
